% =========================================================
% INTRODUCCIÓN
% El encabezado (\chapter*) está en main.tex.
% Referencias en documents/preliminares/preliminares.bib.
% =========================================================

En diciembre de 2025, el 63.8~\% de la población de 18 años o más en México opinó que es inseguro vivir en su ciudad. Esta cifra no es solo un número; representa a millones de personas replanteando sus rutas, evitando salir de noche o caminando con miedo \cite{ENSUdic2025}. Para quienes necesitan moverse en una zona conocida, puede resultar sencillo escoger entre una calle u otra para llegar a un punto de destino; sin embargo, carecer de información sobre el entorno impide planear una ruta segura, dejando al usuario vulnerable ante riesgos derivados de decisiones improvisadas.

La solución estándar que se propone en trabajos anteriores es revisar los datos históricos de delitos y los informes oficiales de incidencias proporcionados por el gobierno para plasmarlos como un indicador de riesgo; sin embargo, este enfoque es estadísticamente inválido en el contexto mexicano. Según la Encuesta Nacional de Victimización y Percepción sobre Seguridad Pública (ENVIPE) 2024, la \enquote{cifra negra} (tasa de criminalidad no declarada) en México se sitúa aproximadamente en 92.9~\% \cite{ENVIPE2024}. Dado que los registros policiales capturan menos del 8~\% de la incidencia delictiva real, una aplicación de navegación basada exclusivamente en ese historial resulta insuficiente para alertar sobre la gran mayoría de los riesgos, generando una falsa sensación de seguridad al usuario. Lo anterior evidencia un sesgo importante en la información que se obtiene de los registros oficiales, por lo que es imperativo encontrar datos alternativos que permitan construir un modelo preciso y confiable.
